\section{Marco Referencial (Pila Tecnológica del Sistema)}

El marco referencial o tecnológico sistematiza las herramientas, plataformas, lenguajes de programación, marcos de trabajo (frameworks) y servicios de infraestructura informática seleccionados para sustentar el diseño de la arquitectura del sistema de gestión y trazabilidad. A diferencia del marco conceptual ---orientado a los principios y modelos teóricos abstractos---, este apartado especifica la viabilidad técnica y las características operativas de la pila tecnológica (stack) propuesta, orientada a respaldar una articulación eficiente entre la persistencia relacional, la lógica de negocio y la integración con modelos de inteligencia artificial.

\subsection{Entorno de Desarrollo y Lenguajes de Programación}

\begin{itemize}
    \item \textbf{PHP \cite{php2024}:} Lenguaje de programación de código abierto del lado del servidor ampliamente utilizado en aplicaciones web corporativas y sistemas transaccionales. Su compatibilidad nativa con sistemas gestores de bases de datos relacionales, su soporte robusto para Programación Orientada a Objetos (POO) y su amplia adopción en arquitecturas de soporte técnico (como se evidencia en los antecedentes de Mosquera Perea et al., \cite{mosquera2025}, y Velasco Castro, \cite{velasco2023}) lo constituyen como el núcleo para la implementación de la lógica del servidor.
    \item \textbf{JavaScript / TypeScript \cite{ecma2023}:} Lenguajes de programación estándar para la ejecución del lado del cliente en navegadores web. Permiten gestionar la asincronía en las solicitudes de estado, dinamizar la carga de interfaces reactivas y procesar la interacción conversacional del usuario en tiempo real sin recargar la página completa.
    \item \textbf{HTML5 y CSS3 \cite{w3c2023}:} Estándares fundamentales para la estructuración semántica del contenido digital y la definición de hojas de estilo en cascada. Garantizan la correcta presentación visual, la accesibilidad de la información de las órdenes de servicio y la compatibilidad multiplataforma en dispositivos de escritorio y móviles.
\end{itemize}

\subsection{Marcos de Trabajo (Frameworks) y Librerías de Interfaz}

\begin{itemize}
    \item \textbf{Laravel \cite{otwell2024}:} Framework de desarrollo web estructurado bajo el patrón arquitectónico Modelo-Vista-Controlador (MVC). Laravel provee mecanismos nativos esenciales para el proyecto: enrutamiento estructurado, motor de mapeo objeto-relacional (ORM Eloquent) para la gestión transparente de consultas a la base de datos, sistema de autenticación y autorización robusto, middleware para la protección contra vulnerabilidades comunes (como inyección SQL y CSRF), y facilidades para la construcción de APIs RESTful que interconecten el backend con servicios externos.
    \item \textbf{Bootstrap 5 \cite{otto2023}:} Framework de diseño frontend enfocado en la creación de interfaces responsivas bajo el paradigma mobile-first. Permite estructurar tableros de control (dashboards) intuitivos y adaptables a pantallas de computadoras, tabletas y teléfonos móviles, facilitando que tanto el personal técnico en taller como los clientes desde sus dispositivos personales puedan interactuar con la plataforma sin barreras visuales.
\end{itemize}

\subsection{Motor de Base de Datos y Persistencia Relacional}

\begin{itemize}
    \item \textbf{MySQL 8.0 \cite{oracle2024}:} Sistema de Gestión de Bases de Datos Relacionales (RDBMS) de alto rendimiento, confiabilidad y soporte para transacciones ACID (Atomicidad, Consistencia, Aislamiento y Durabilidad). Se selecciona para el diseño del esquema entidad-relación del sistema, orientado a respaldar la integridad referencial entre clientes, dispositivos, diagnósticos, órdenes de trabajo, técnicos asignados y marcas de tiempo (timestamps). Asimismo, contempla implementar restricciones de borrado lógico (soft deletes) a nivel de ORM y un esquema de auditoría mediante tablas de registro (audit\_logs) y permisos restrictivos de usuario DBMS (revocación de permisos DELETE), orientados a fortalecer la integridad del sistema y preservar un historial auditable reduciendo el riesgo de alteraciones no autorizadas.
    \item \textbf{MySQL Workbench:} Herramienta visual de modelado, administración y diseño de bases de datos que facilita la creación de esquemas físicos, diagramas entidad-relación y ejecución controlada de scripts de inicialización y migración.
\end{itemize}

\subsection{Herramientas de Modelado de Software y Documentación}

\begin{itemize}
    \item \textbf{Enterprise Architect / StarUML \cite{omg2017}:} Herramientas especializadas de modelado basadas en el estándar UML 2.5. Se emplearán para elaborar los cinco diagramas formales de ingeniería de software contemplados en la arquitectura del sistema: Diagramas de Casos de Uso (definición de roles y límites del sistema), Diagramas de Actividades (flujos de trabajo del ciclo de reparación), Diagramas de Clases (modelo de dominio y persistencia), Diagramas de Secuencia (intercambio temporal de mensajes) y Diagramas de Despliegue (topología física de hardware, servidor web Laravel, RDBMS MySQL y API de IA).
\end{itemize}

\subsection{Pila Tecnológica de Inteligencia Artificial Conversacional}

\begin{itemize}
    \item \textbf{Modelo de Lenguaje de Gran Tamaño en la Nube (Google Gemini API REST):} Se establece como alternativa tecnológica de referencia la plataforma Google Gemini \cite{googleai2026}, proyectada para su integración mediante API REST con respuestas estructuradas en formato JSON, cuya selección y versión definitiva de modelo se evaluará en función de su disponibilidad, estabilidad de servicio, latencia de procesamiento y rendimiento en español.
    \item \textbf{Función de Interpretación y Redacción del LLM:} En el diseño propuesto, el modelo de lenguaje actuará estrictamente como motor de comprensión semántica e interfaz conversacional en español. No poseerá permisos de ejecución de código ni acceso directo a las tablas del sistema.
    \item \textbf{Mecanismo de Desacoplamiento y Control de Acceso (API REST HTTPS Intermediaria):} Capa de software intermedia proyectada en el backend de Laravel para validar el token verificado de la orden de servicio, realizar lecturas seguras de solo lectura en MySQL y proveer los datos necesarios al LLM dentro de un entorno controlado mediante ingeniería de instrucciones (system prompts).
\end{itemize}